\chapter{線形写像}
    \section{表記の規則}
        \subsection{一般のベクトルを要素とする配列と行列の積}
            $\bm{v}_1,\dotsc,\bm{v}_n$をベクトルとする($\realNumbers^n$や$\complexNumbers^n$の要素とは限らないことに注意)。
            これと$m\times n$行列$A=[a_{ij}]_{m\times n}$, $n\times m$行列$B=[b_{ij}]_{n\times m}$に対して次のように演算を定義する。
            \[
                A
                \begin{bmatrix}
                    \bm{v}_1 \\
                    \vdots \\
                    \bm{v}_n
                \end{bmatrix} \coloneqq
                \begin{bmatrix}
                    \sum_{j=1}^n a_{1j}\bm{v}_j \\
                    \vdots \\
                    \sum_{j=1}^n a_{mj}\bm{v}_j
                \end{bmatrix},
                \qquad
                [\bm{v}_1,\dotsc,\bm{v}_n]B \coloneqq
                \left[\sum_{i=1}^n b_{i1}\bm{v}_i, \dotsc, \sum_{i=1}^n b_{im}\bm{v}_i \right]
            \]
            これは$\bm{v}_1,\dotsc,\bm{v}_n$を形式的にスカラーと見做して普通の行列同士の積を考えたものである。
        \subsection{一般のベクトルを要素とする行列同士の積}
            ベクトル$\bm{a}_{ij}\;((i,j)\in\{1:l\}\times\{l:m\})$と$\bm{b}_{ij}\;((i,j)\in\{1:m\}\times\{l:n\})$に対して、
            これらを要素とする行列$A\coloneqq[\bm{a}_{ij}]_{l\times m},\;B\coloneqq[\bm{b}_{ij}]_{m\times n}$の積$C$を
            $C\coloneqq[c_{ij}]_{m\times n},\;c_{ij} \coloneqq \sum_{k=1}^l \inprod{\bm{a}_{ik}}{\bm{b}_{kj}}$で定義する。
            これは$\bm{a}_{ij},\bm{b}_{ij}$を形式的にスカラーと見做して普通の行列同士の積を考えたものである。
    \section{\texorpdfstring{$\dim(V)=\dim(W),\;f:V\to W$}{dim(V)=dim(W), f:V→W}であるとき、\texorpdfstring{$f$}{f}が単射であることと全射であることは同値。}
        \begin{shadebox}
            $f$を$V\rightarrow W$への線形写像とする。$\dim(V)=\dim(W)$であるとき、$f$が単射であることと全射であることは同値。
        \end{shadebox}
        \begin{proof}
            \quad\par
            \noindent
            $n=\dim(V)=\dim(W)$とする。\\
            ($\Rightarrow$)
            \par
            $W$の任意のベクトルは$W$の基底の一次結合で表されるから、全射性を言うためには$V$の中から適当な$n$個のベクトルを選んで$f$で送ったものの集合で$W$の基底を構成できることを示せば十分である。
            $V$の基底を$\{\bm{v}_1,\bm{v}_2,\dots,\bm{v}_n\}$として$\{f(\bm{v}_1),f(\bm{v}_2),\dots,f(\bm{v}_n)\}$が$W$の基底を成すことを示す。
            \[c_1f(\bm{v}_1)+c_2f(\bm{v}_2)+\dotsb+c_nf(\bm{v}_n) = 0,\;\;\;c_1,c_2,\dots c_n\in\field\]
            とおくと$f$の線形性から
            \[f(c_1\bm{v}_1+c_2\bm{v}_2+\dotsb+c_n\bm{v}_n) = 0\]
            $f$は単射であったから$\Ker{f}=\{\bm{0}\}$なので
            \[c_1\bm{v}_1+c_2\bm{v}_2+\dotsb+c_n\bm{v}_n = 0\]
            $\{\bm{v}_1,\bm{v}_2,\dots,\bm{v}_n\}$は$V$の基底なので
            \[c_1=c_2=\dotsb=c_n = 0\]
            よって$W$の$n=\dim(W)$個のベクトル$\{f(\bm{v}_1),f(\bm{v}_2),\dots,f(\bm{v}_n)\}$が一次独立なのでこれは$W$の基底である。
            \newline\par
            \noindent
            ($\Leftarrow$)
            \par
            次元定理から
            \[\dim(\Img{f}) + \dim(\Ker{f}) = \dim(V) = \dim(W) = n\]
            $f$は全射であったから$\Img{f}=W$であるので$\dim(\Img{f}) = n$。よって上式より$\dim(\Ker{f}) = 0$なので$f$は単射である。\\
            (別証)
            \par
            $f$が全射であるから、$W$の任意の基底$\{\bm{w}_1,\dotsc,\bm{w}_n\}$に対して適当な$V$のベクトル$\bm{v}_1,\dotsc,\bm{v}_n$が存在して
            $f(\bm{v}_i)=\bm{w}_i\;(i=1,\dotsc,n)$を満たす。
            ここで$c_1\bm{v}_1 + \dotsb + c_n\bm{v}_n = \bm{0}$とおくと、$f$は線形写像なので$f(\bm{0})=\bm{0}$だから$\bm{0} = f(c_1\bm{v}_1 + \dotsb + c_n\bm{v}_n) = c_1f(\bm{v_1}) + \dotsb + c_nf(\bm{v_n}) = c_1\bm{w}_1 + \dotsb + c_n\bm{w}_n$となる。
            $\{\bm{w}_1,\dotsc,\bm{w}_n\}$は$W$の基底であったから$c_1=\dotsb=c_n=0$である。
            よって$\{\bm{v}_1,\dotsc,\bm{v}_n\}$は一次独立であり、$\dim{V}=n$よりこれらは$V$の基底である。
            \par
            $\bm{x}\in V$に対して$f(\bm{x})=\bm{0}$とおく。
            $\bm{x}$を先程の$V$の基底で展開して$\bm{x}=c_1\bm{v}_1+\dotsb+c_n\bm{v}_n$とすると、
            \[\bm{0} = f(\bm{x}) = f(c_1\bm{v}_1+\dotsb+c_n\bm{v}_n) = c_1\bm{w}_1 + \dotsb + c_n\bm{w}_n\]
            より$c_1=\dotsb=c_n=0$となるから$\bm{x}=\bm{0}$である。
            すなわち$\Ker{f}=\{\bm{0}\}$なので$f$は単射である。
        \end{proof}
    \section{\texorpdfstring{$\dim{V}=\dim{W}$}{dim(V)=dim(W)}ならば全単射線形写像\texorpdfstring{$f:V\to W$}{f:V→W}が存在する}
        \begin{proof}
            \quad\par
            $n=\dim{V}=\dim{W}$とし、$\{\bm{v}_1,\dotsc,\bm{v}_n\}$を$V$の基底、$\{\bm{w}_1,\dotsc,\bm{w}_n\}$を$W$の基底とする。
            写像$f:V\to W$を$f(\bm{v}_i)=\bm{w}_i\;(i=1,\dotsc,n),\;^\forall\bm{x},\bm{y}\in V,\alpha,\beta\in\realNumbers,f(\alpha\bm{x}+\beta\bm{y})=\alpha f(\bm{x})+\beta f(\bm{y})$で定めると$f$は線形写像である。
            さらに単射であることが次のようにして確かめられる。
            \par
            $\bm{x},\bm{y}\in V$に対して$f(\bm{x})=f(\bm{y})$であるとする。
            $\bm{x},\bm{y}$を基底で展開して$\bm{x}=\sum_{i=1}^n x_i\bm{v}_i, \bm{y}=\sum_{i=1}^n y_i\bm{v}_i$とすると$\bm{x}=\sum_{i=1}^n x_i\bm{w}_i = \bm{y}=\sum_{i=1}^n y_i\bm{w}_i$すなわち$\bm{x}=\sum_{i=1}^n (x_i-y_i)\bm{w}_i = \bm{0}$であるが、$\{\bm{w}_1,\dotsc,\bm{w}_n\}$は$W$の基底であり一次独立だから$x_i=y_i\;(i=1,\dotsc,n)$であるので$\bm{x}=\bm{y}$である。直前の定理より$f$は全射でもあるから全単射である。
        \end{proof}
    \section{不変部分空間と表現行列}
        \label{不変部分空間と表現行列}
        \begin{shadebox}
            $A\in\field^{n\times n},\quad W_1,\dotsc,W_m\in\field^n : A$の不変部分空間,\quad$\field^n = W_1\oplus\dotsb\oplus W_m,\quad d_i \coloneqq \dim{W_i}$\\
            $\{\bm{w}_{i1},\dotsc,\bm{w}_{id_i}\} : W_i$の基底,\quad$\Phi_i \coloneqq [\bm{w}_{i1},\dotsc,\bm{w}_{id_i}],\quad \Phi \coloneqq [\Phi_1,\dotsc,\Phi_m]$\\
            以上の設定の下で、次式を満たす適当な正方行列$A_i\in\field^{d_i\times d_i}$が存在する。
            \[A = \Phi\;\diag{A_1,\dotsc,A_m}\Phi^{-1}\]
            $A_i$は次式で得られる。
            \[A_i = \left({\Phi_i}^\top\Phi_i\right)^{-1}{\Phi_i}^\top A\Phi_i\]
            特に$W_i$の基底を正規直交基底にしておけば、より簡単に次式で得られる。
            \[A_i = {\Phi_i}^\top A\Phi_i\]
        \end{shadebox}
        \begin{proof}
            \quad\par
            $A = \Phi\;\diag{A_1,\dotsc,A_m}\Phi^{-1}$を満たす$A_i$が存在することは\cite{システム線形代数}定理1.20を見ればわかる。
            $\Phi$は$\field$から来る入力ベクトルを$A$の不変部分空間の成分に分解する変換器になっている。
            \par
            $A_i$の構成法を示す。一般の場合を示すのは別に難しくないが、式が煩雑になって無意味に難解になるので、具体的な場合を示す。
            例えば$n=4, m=2, d_1=d_2=2$としてみる。この場合の証明が理解できれば一般の証明は楽勝である。
            まず
            \[A = \Phi\;\diag{A_1,\dotsc,A_m}\Phi^{-1} = [\bm{w}_1,\bm{w}_2,\bm{w}_3,\bm{w}_4]\diag{A_1,A_2}[\bm{w}_1,\bm{w}_2,\bm{w}_3,\bm{w}_4]^{-1}\]
            とおく。まず$A_1$を求めてみる。定理\ref{逆行列にその構成縦ベクトルを掛ける}からヒントを得て両辺に$[\bm{w}_1,\bm{w}_2]$を右から掛けて
            \begin{align*}
                A[\bm{w}_1,\bm{w}_2] &= [\bm{w}_1,\bm{w}_2,\bm{w}_3,\bm{w}_4]\diag{A_1,A_2}
                \left[
                    \begin{array}{c}
                        I_2\\
                        O_2
                    \end{array}
                \right] = [\bm{w}_1,\bm{w}_2,\bm{w}_3,\bm{w}_4]
                \left[
                    \begin{array}{c}
                        A_1\\
                        O_2
                    \end{array}
                \right]\\
                &= [\bm{w}_1,\bm{w}_2]A_1
            \end{align*}
            $[\bm{w}_1,\bm{w}_2]$は列フルランクであるから定理\ref{列フルランク行列を掛けあわせて正則行列を作る}より$[\bm{w}_1,\bm{w}_2]^\top[\bm{w}_1,\bm{w}_2]$は逆行列を持つ。
            そこで上式に左から$[\bm{w}_1,\bm{w}_2]^\top$を掛けて
            \begin{align*}
                [\bm{w}_1,\bm{w}_2]^\top A[\bm{w}_1,\bm{w}_2] &= [\bm{w}_1,\bm{w}_2]^\top[\bm{w}_1,\bm{w}_2]A_1\\
                \therefore\;A_1 &= ([\bm{w}_1,\bm{w}_2]^\top[\bm{w}_1,\bm{w}_2])^{-1}[\bm{w}_1,\bm{w}_2]^\top A[\bm{w}_1,\bm{w}_2]
            \end{align*}
            特に、予め$\{\bm{w}_1,\bm{w}_2\}$を正規直交基底にしておけば$[\bm{w}_1,\bm{w}_2]^\top[\bm{w}_1,\bm{w}_2] = I_2$なので$A_1$は次式で得られる。
            \[A_1 = [\bm{w}_1,\bm{w}_2]^\top A[\bm{w}_1,\bm{w}_2]\]
            同様にして$A_2$も求められる。一般の場合も同じ要領でやれば良い。
        \end{proof}
    \subsection{正規直交基底間での座標変換は直交変換}
        \begin{shadebox}
            $n$次元線形空間$V$の2つの正規直交基底$\{\bm{u}_1,\dots,\bm{u}_n\},\;\{\bm{v}_1,\dotsc,\bm{v}_n\}$に対して、$\{\bm{u}_i\}$座標系から$\{\bm{v}_i\}$座標系への座標変換は直交変換である。
            すなわち座標変換行列$P \coloneqq [\inprod{\bm{u}_i}{\bm{u}_j}]_{n\times n}$は直交行列である。
        \end{shadebox}
        \begin{proof}
            \[ \bm{v}_i = \sum_{k=1}^n \inprod{\bm{v}_i}{\bm{u}_k}\bm{u}_k = \sum_{k=1}^n P[i][k]\bm{u}_k \]
            だから
            \begin{align*}
                \delta_{ij} &= \inprod{\bm{v}_i}{\bm{v}_j} = \inprod{\sum_{k=1}^n P[i][k]\bm{u}_k}{\sum_{l=1}^n P[j][l]\bm{u}_l} = \sum_{k=1}^n\sum_{l=1}^n \inprod{P[i][k]\bm{u}_k}{P[j][l]\bm{u}_l}\\
                &= \sum_{k=1}^n\sum_{l=1}^n P[i][k]P[j][l]\delta_{kl} = \sum_{k=1}^n P[i][k]P[j][k] = \sum_{k=1}^n P[i][k]P^\top[k][j] = (PP^\top)[i][j]
            \end{align*}
        \end{proof}